\documentclass[10pt,a4paper]{amsart}
\usepackage[margin=19mm]{geometry}
\usepackage{amsmath,amssymb,mathtools}
\usepackage[scr=rsfs,cal=euler]{mathalfa}
\usepackage{xfrac}
\usepackage[unicode,hidelinks,bookmarks=false]{hyperref}
\newcommand{\ie}{i.e.\ }
\newcommand{\cf}{cf.\ }
\newcommand{\resp}{respectively\ }
\newcommand{\Subsection}{\subsection}
\providecommand{\nobreakdash}{\nobreak\hbox{-}}
\setlength{\emergencystretch}{2em}
\multlinegap0pt
\usepackage{color}
\usepackage{multicol}
\usepackage{mathtools}
\usepackage{tikz,tikz-cd,tcolorbox}
\usepackage{setspace,mathrsfs}
\newtheorem{theorem}{Theorem}[section]
\newtheorem{lemma}[theorem]{Lemma}
\newtheorem{proposition}[theorem]{Proposition}
\newtheorem{corollary}[theorem]{Corollary}
\usepackage{amsaddr}
\newtheorem{definition}[theorem]{Definition}
\newtheorem{remark}[theorem]{Remark}
\newtheorem{example}[theorem]{Example}
\newtheorem{notation}[theorem]{Notation}
\newcommand{\Max}[1]{{\color{red}Max:#1}}
\newcommand{\Joseph}[1]{{\color{blue}Joseph:#1}}
\begin{document}
\title[Duality for partial orthomodular lattices]{Duality for partial orthomodular lattices}
\author{Joseph McDonald}
\date{}
\maketitle
\noindent\emph{Source and licence.} Duality for partial orthomodular lattices. Version arXiv:2608.05090v1 (CC BY 4.0). This material is distributed under the Creative Commons Attribution 4.0 International licence, http://creativecommons.org/licenses/by/4.0/, as recorded in the arXiv OAI licence metadata for this exact version and shown on the abstract page https://arxiv.org/abs/2608.05090v1. Full rights notice: licenses/arxiv-2608.05090-CC-BY-4.0.txt.\par\smallskip
\noindent\textit{Editorial scope.} Counted unit: the original \texttt{theorem} environment \texttt{homeomorphism} at source lines 127--129 together with its complete proof at lines 130--140; the delivered document keeps source lines 47--141 so that every referenced label and every citation resolves inside it. Native editable source is the paper's own concatenated TeX; the AMS wrapper is editorial formatting only. No publisher class, upstream script or unrelated artwork is redistributed.
\section{Topological representation and duality}
Recall that an OL is a bounded lattice $\langle A;\wedge,\vee,0,1\rangle$ equipped with an order-inverting period two complementation $^\perp\colon A\to A$, known as an \emph{orthocomplementation}. An OML is an OL satisfying $a\leq b\Rightarrow b=a\vee(a^{\perp}\wedge b)$. 
\begin{definition}
    A \emph{partial orthomodular lattice} (\emph{pOML}) is an OL such that  $b\leq a$ and $a=b\vee(a\wedge b^{\perp})$ implies $b=a\wedge(a^{\perp}\vee b)$.  
\end{definition}
%\begin{example}\label{example}
%The ortholattice of closed convex cones in $\mathbb{R}^2$ is partial orthomodular under the operations of intersection and orthogonal complementation. This ortholattice is however not orthomodular (see \cite[Example 1]{leem}).   
%\end{example}

   Note that not every OL is a pOML but every OML is a pOML (see \cite[Propositions 2 and 3]{leem}).                                    An \emph{orthoframe} is a pair $\langle X;\perp\rangle$ such that $X$ is a set and $\perp\subseteq X^2$ is irreflexive and symmetric. For any orthoframe $X$ and $U\subseteq X$, let \[U^{\perp}:=\{x\in X:x\perp y\hspace{.1cm}\text{for all}\hspace{.1cm}y\in U\}.\] In \cite{goldblatt2} a subset $U\subseteq X$ is defined as $\perp$-\emph{closed} relative to a subset $V\subseteq X$ if for all $x\in V$, if 
    $x\not\in U$, there exists $y\in V$ such that $y\perp U$ and $x\not\perp y$. We denote the sets that are $\perp$-closed relative to the entire space $X$ by $\mathcal{B}(X)$. Note that $\mathcal{B}(X)$ is precisely those subsets that are biorthogonally closed in the sense that $\mathcal{B}(X)=\{U\subseteq X:U=U^{\perp\perp}\}$.     
\begin{definition}
    A $c$\emph{-frame} is a triple $\langle X;\perp,\Omega(X)\rangle$ such that $\langle X;\perp\rangle$ is an orthoframe and $\Omega(X)\subseteq\mathcal{B}(X)$ is closed under $\cap$, $^{\perp}$, and satisfies the following condition: if $V\subseteq U$ and $U^{\perp}$ is $\perp$-closed in $V^{\perp}$, then $V$ is $\perp$-closed in $U$.
\end{definition}
\begin{lemma}\label{lemma 2.3}
    If $X$ is a c-frame, then $U=V\cap(V^{\perp}\sqcup U)$ implies that $U$ is $\perp$-closed in $V$ for all $U,V\in\Omega(X)$ where $U\sqcup V:=(U\cup V)^{\perp\perp}$.  
\end{lemma}
\begin{proof}
    Take $x\in V$ with $x\not\in U$. Since $U=V\cap(V^{\perp}\sqcup U)$, we have $x\not\in V\cap(V^{\perp}\sqcup U)$. Moreover, since $x\in V$, it follows that $x\not\in V^{\perp}\sqcup U=(V\cap U^{\perp})^{\perp}$ and thus there exists some $y\in V\cap U^{\perp}$ such that $x\not\perp y$. Hence $y\in V$ and $y\in U^{\perp}$ which implies $y\perp U$. Therefore $U$ is $\perp$-closed in $V$.
    \end{proof}
\begin{lemma}\label{c-frame to poml}
    If $X$ is a c-frame, then $\langle\Omega(X);\cap,\sqcup,^{\perp},\emptyset,X\rangle$ is a pOML. 
\end{lemma}
    \begin{proof}
Since $\Omega(X)\subseteq\mathcal{B}(X)$, it follows that $\langle\Omega(X);\cap,\sqcup,^{\perp},\emptyset,X\rangle$ forms an OL. It remains to verify that it is a pOML. Assume $U,V\in\Omega(X)$ with $V\subseteq U$ and $U=V\sqcup(U\cap V^{\perp})$. We first show $U^{\perp}$ is $\perp$-closed in $V^{\perp}$. Since $^{\perp}$ is an orthocomplementation, De Morgan's identities give $U^{\perp}=V^{\perp}\cap(V\sqcup U^{\perp})$. Then Lemma \ref{lemma 2.3} implies $U^{\perp}$ is $\perp$-closed in $V^{\perp}$. Since $X$ is a c-frame, it follows that $V$ is $\perp$-closed in $U$. To show $V=U\cap(U^{\perp}\sqcup V)$, we demonstrate the $U\cap(U^{\perp}\sqcup V)\subseteq V$ inclusion as the other inclusion is trivial. Hence assume $x\not\in V$. Since $V$ is $\perp$-closed in $U$, there exists $y\in U$ such that $x\not\perp y$ and $y\perp V$. Moreover, we have $y\in U\cap V^{\perp}$. Now suppose for the sake of contradiction that $x\in U\cap(U^{\perp}\sqcup V)$ i.e., $x\in U\cap(U\cap V^{\perp})^{\perp}$. Then $x\perp z$ for every $z\in U\cap V^{\perp}$ but this contradicts our hypothesis that $x\not\perp y$ for some $y\in U\cap V^{\perp}$. Hence $x\not\in U\cap(U^{\perp}\sqcup V)$ so $U\cap(U^{\perp}\sqcup V)\subseteq V$.  
\end{proof} 
\begin{definition}
    Let $A$ be a pOML. Then define $\mathcal{F}(A)=\langle X_A;\perp_A,\beta_A,\tau_A\rangle$ by: 
    \begin{enumerate}
        \item $X_A$ is the collection of all non-empty proper filters of $A$; 
        \item $x\perp_A y$ iff there exists $a\in A$ such that $a\in x$ and $a^{\perp}\in y$; 
        \item $\beta_A=\{\phi(a):a\in A\}$ where $\phi(a)=\{x\in X_A:a\in x\}$; 
        \item $\tau_A$ is the topology on $X_A$ generated by the basis $\beta_A$. 
    \end{enumerate}
\end{definition}
For a topological space $X$, let $\mathcal{C}(X)$ be the compact subsets of $X$, let $\mathcal{O}(X)$ be the open subsets of $X$, and let $\mathcal{K}(X):=\mathcal{C}(X)\cap\mathcal{O}(X)$. Recall that $X$ is a \emph{spectral space} if $X$ is compact, $T_0$, coherent, and sober (see \cite[Section 3.2]{mcdonald}). Since every spectral space $X$ is $T_0$, one can define a specialization order by $x\leqslant y$ iff $x\in U$ implies $y\in U$ for all $U\in\mathcal{O}(X)$, that is a partial order.  
\begin{lemma}\label{spectral space}
     $\langle X_A;\tau_A\rangle$ is a spectral space whose specialization order is $\subseteq$.  
\end{lemma}
\begin{proof}
Since every pOML is an OL, the result is an immediate application of \cite[Proposition 3.4.1]{mcdonald}.  
\end{proof}
\begin{remark}
    It was shown in \cite[Proposition 3]{goldblatt1} that for any OL $A$, the Stone space $\langle X_A;\tau_A^*\rangle$ is compact where $\tau_A^*\subseteq\wp(X_A)$ is the topology generated by the subbasis $\bigcup_{a\in A}\{\phi(a),\hspace{.1cm}X_A\setminus\phi(a)\}$. By \cite[Proposition 3.3.1]{mcdonald}, the compactness of $\langle X_A;\tau_A^*\rangle$ requires the Axiom of Choice. As the proof of compactness of $\langle X_A;\tau_A\rangle$ follows by \cite[Proposition 3.4.1]{mcdonald}, it obtains in ZF alone.     
\end{remark}
\begin{lemma}\label{lemma 2.8}
    In any OL, we have $b\wedge(b^{\perp}\vee a)\leq a$ iff $\phi(a)$ is $\perp$-closed in $\phi(b)$. 
\end{lemma}
\begin{proof}
    This is an application of \cite[Lemma 6.3]{goldblatt2}, \cite[Theorem 3.7]{goldblatt2}, and the fact that OLs provide sound and complete algebraic models for orthologic.      
\end{proof}
\begin{lemma}\label{poml to c-frame}
    $\langle X_A;\perp_A,\beta_A\rangle$ is a c-frame. 
\end{lemma}
\begin{proof}
    The proof that $\langle X_A;\perp_A\rangle$ is a orthoframe follows from \cite[Proposition 2]{goldblatt1}. It suffices to show that for all $a,b
    \in A$, if $\phi(b)\subseteq\phi(a)$ and $\phi(a)^{\perp}$ is $\perp$-closed in $\phi(b)^{\perp}$, then $\phi(b)$ is $\perp$-closed in $\phi(a)$. First note that the definition of $\phi$ applied to our hypothesis that $\phi(b)\subseteq\psi(a)$ gives $b\leq a$. Applying Lemma \ref{lemma 2.8} to our hypothesis that $\phi(a)^{\perp}$ is $\perp$-closed in $\phi(b)^{\perp}$ yields $b^{\perp}\wedge(b^{\perp\perp}\vee a^{\perp})\leq a^{\perp}$. Since $^{\perp}$ is an order-inverting involution, $a^{\perp\perp}=a\leq(b^{\perp}\wedge(b^{\perp\perp}\vee a^{\perp}))^{\perp}$. De Morgan's identities and the fact that $^{\perp}$ is an involution then guarantees that  \[(b^{\perp}\wedge(b^{\perp\perp}\vee a^{\perp}))^{\perp}=b\vee(b^{\perp}\wedge a)\] so $a\leq b\vee(b^{\perp}\wedge a)$. Clearly $b^{\perp}\wedge a\leq a$ which together with the fact that $b\leq a$ and the anti-symmetry of $\leq$ yields $a=b\vee(b^{\perp}\wedge a)$. Since $A$ is a pOML, we have $b=a\wedge(a^{\perp}\vee b)$ so $a\wedge(a^{\perp}\vee b)\leq b$. Lemma \ref{lemma 2.8} implies $\phi(b)$ is $\perp$-closed in $\phi(a)$.  
\end{proof}
For a $c$-frame $X$ equipped with a topology, we define \[\mathcal{K}\Omega(X):=\mathcal{K}(X)\cap\Omega(X).\] 
\begin{theorem}\label{isomorphism}
    Every pOML $A$ is isomorphic to $\mathcal{K}\Omega(\mathcal{F}(A))$.  
\end{theorem}
\begin{proof}
    The map $\phi(a)=\{x\in X_A:a\in x\}$ provides an isomorphism. The proof is a simple application of \cite[Theorem 3.4.2]{mcdonald}, however we outline the proof that $\phi$ is a bijection. For any $c\in A$, let ${\uparrow}c=\{d\in A:c\leq d\}$. Clearly ${\uparrow}c$ is a proper filter whenever $c\not=0$. For injectivity, assume $a\not\leq b$. Then $a\in{\uparrow}a$ but $b\not\in{\uparrow}a$ so ${\uparrow}a\in\phi(a)$ but ${\uparrow}a\not\in\phi(b)$ and hence $\phi(a)\not\subseteq\phi(b)$. For surjectivity, suppose $U\in\mathcal{K}\Omega(\mathcal{F}(A))$ so that $U=\bigcup^n_{i=1}\phi(a_i)$ for $a_1,\dots a_n\in A$. Since $U\in\Omega(\mathcal{F}(A))$, it can be easily shown that $U=U^{\perp\perp}$ and hence we obtain \[U=U^{\perp\perp}=\biggl(\bigcup_{i=1}^n\phi(a_i)\biggl)^{\perp\perp}=\bigvee^n_{i=1}\phi(a_i)=\phi\biggl(\bigvee_{i=1}^na_i\biggl).\] Therefore $U$ is in the image of $\phi$ and hence $\phi$ is surjective. 
\end{proof}
 Below are a restriction to the upper Vietoris orthospaces (the choice-free duals of the OLs) introduced in \cite{mcdonald}. For the following definition, let \[\mathcal{K}\Omega_X(x):=\{U\in\mathcal{K}\Omega(X):x\in U\}.\]  
\begin{definition}\label{uvc-space}
    An \emph{upper Vietoris c-space} (\emph{UVC-space}) is a c-frame $X$ equipped with a $T_0$-topology $\tau\subseteq\wp(X)$ satisfying the following conditions: 
    \begin{enumerate}
        \item $\mathcal{K}\Omega(X)$ is closed under $\cap$ and $^{\perp}$ and forms a basis for $X$;
        \item if $x\perp y$, there exists $U\in\mathcal{K}\Omega(X)$ such that $x\in U$ and $y\in U^{\perp}$; 
        \item every proper filter in $\mathcal{K}\Omega(X)$ is of the form $\mathcal{K}\Omega_X(x)$ for some $x\in X$. 
    \end{enumerate}
\end{definition}
\begin{lemma}\label{functors}
    If $X$ is a UVC-space, then $\mathcal{G}(X)=\langle\mathcal{K}\Omega(X);\cap,\sqcup,^{\perp},\emptyset,X\rangle$ is a pOML. If $A$ is a pOML, then $\mathcal{F}(A)=\langle X_A;\perp_A,\beta_A,\tau_A\rangle$ is a UVC-space. 
\end{lemma}
\begin{proof}
    The first part follows by Lemma \ref{c-frame to poml} and \cite[Lemma 4.1.2]{mcdonald}. For the second part, by Lemma \ref{spectral space} and Lemma \ref{poml to c-frame}, it suffices to show that conditions 1-3 of Definition \ref{uvc-space} are satisfied. For condition 1, it follows by Theorem \ref{isomorphism} that $\phi(a)\cap\phi(b)=\phi(a\wedge b)\in\mathcal{K}\Omega(\mathcal{F}(A))$ and $\phi(a)^{\perp}=\phi(a^{\perp})\in\mathcal{K}\Omega(\mathcal{F}(A))$. Moreover, the definition of $\beta_A$ and $\tau_A$ guarantees that sets of the form $\phi(a)$ for $a\in A$ are a basis for $\mathcal{F}(A)$. For condition 2, assume $x\perp_Ay$. Then there exists some $a\in A$ such that $a\in x$ and $a^{\perp}\in y$. Therefore we have $x\in\phi(a)$ and $y\in\phi(a^{\perp})=\phi(a)^{\perp}$. For condition 3, take any proper filter $F$ in $\mathcal{K}\Omega(\mathcal{F}(A))$. Theorem \ref{isomorphism} gives that $x=\{a\in A:\phi(a)\in F\}$ is a proper filter in $A$. Therefore $x\in\mathcal{F}(A)$ and $\mathcal{K}\Omega_{\mathcal{F}(A)}(x)=F$, as desired.    
\end{proof}

\begin{theorem}\label{homeomorphism}
    Every UVC-space $X$ is homeomorphic and relationally isomorphic to $\mathcal{F}(\mathcal{G}(X))$. Moreover every UVC-space is a spectral space.  
\end{theorem}

\begin{proof}
    We show that the map $\psi(x)=\{U\in\mathcal{K}\Omega(X):x\in U\}$ exhibits the desired homeomorphism and relational isomorphism. It is an easy exercise to verify that $\psi(x)$ is a proper filter in $\mathcal{G}(X)$ for any $x\in X$ so that $\psi$ is well-defined. To see that $\psi$ is injective, note that since $\mathcal{K}\Omega(X)$ forms a basis for $X$, it follows that if $x\not\leqslant y$, there exists $U\in\mathcal{K}\Omega(X)$ such that $x\in U$ and $y\not\in Y$. The definition of $\psi$ then gives $U\in\psi(x)$ and $U\not\in\psi(y)$ so $\psi(x)\not\leqslant\psi(y)$. For surjectivity, take any $F\in\mathcal{F}(\mathcal{G}(X))$. By Definition \ref{uvc-space}(3), there exists some $x\in X$ such that $\mathcal{K}\Omega_X(x)=F$ and hence $\psi(x)=F$. To see $\psi$ is continuous, take any basic open $\phi(U)$ in $\mathcal{F}(\mathcal{G}(X))$, then 
    \begin{align*}
        \psi^{-1}[\phi(U)]&=\{x\in X:\mathcal{K}\Omega_X(x)\in\phi(U)\}\\&=\{x\in X:U\in\mathcal{K}\Omega_X(x)\}\\&=\{x\in X:x\in U\}\\&=U
    \end{align*}
   The continuity of $\psi^{-1}$ is calculated by noting that 
 \[
       \psi[\phi(U)]=\{\mathcal{K}\Omega_X(x):x\in U\}=\{\mathcal{K}\Omega_X(x):U\in\mathcal{K}\Omega_X(x)\}=\phi(U)
\]
 Lastly, it can be shown that $x\perp y$ iff $\psi(x)\perp\psi(y)$ by Definition \ref{uvc-space}(2) and that $x\leqslant y$ iff $\psi(x)\subseteq\psi(y)$ since $x\not\leqslant y$ implies there exists $U\in\mathcal{K}\Omega(X)$ such that $x\in U$ and $y\not\in U$ for any UVC-space. Hence $\psi$ is a relational isomorphism. Thus, it follows by Lemma \ref{spectral space} and Lemma \ref{functors} that $X$ is a spectral space.     
\end{proof}
 

\begin{thebibliography}{99}
\bibitem[goldblatt1]{goldblatt1} Goldblatt, R.I., The Stone space of an ortholattice. Bull. Lond. Math. Soc. \textbf{7}, 45--48 (1975)
\bibitem[goldblatt2]{goldblatt2} Goldblatt, R.I., A semantic analysis of orthologic. J. Phil. Log. \textbf{3}, 19--35 (1974)
\bibitem[leem]{leem} Leemhuis, M., Wolter, D., Özçep, Ö.L.: Rules of Partial Orthomodularity. In: Metcalfe, G., Studer, T., de Queiroz, R. (eds) Logic, Language, Information, and Computation. WoLLIC 2024. Lecture Notes in Computer Science, vol 14672., 108--121, Springer, Cham. (2024)
\bibitem[mcdonald]{mcdonald} McDonald, J., Yamamoto, K.: Choice-free duality for orthocomplemented lattices by means of spectral spaces. \emph{Algebra Universalis}, vol. \textbf{83} (2022)
\end{thebibliography}
\end{document}
